% ── Section I: Introduction ──────────────────────────────────────────
% Length  : ~1 page (double-column IEEE = ~700–900 words)
% Structure:
%   1. Hook       — Vietnam traffic fatality statistics
%   2. Problem    — 2 concrete failure cases from real footage
%   3. Gap        — table / list of what prior systems miss
%   4. Solution   — 1-paragraph plain-English overview of pipeline
%   5. Contributions — C1, C2, C3 explicit bullet points
%   6. Outline    — "The rest of this paper is organized as follows…"
% ─────────────────────────────────────────────────────────────────────

\section{Introduction}
\label{sec:intro}

% [HOOK] Vietnam + SEA traffic context
% TODO: cite WHO or Vietnam Ministry of Transport accident statistics
Vietnam records one of the highest road accident fatality rates in
Southeast Asia, with motorcycles accounting for approximately 94\% of
urban traffic flow~\cite{TODO_VN_traffic_stats}.
%
Automated traffic surveillance systems that can detect accidents
in real time and provide high-quality observations for adaptive signal
control are therefore of practical urgency.

% [PROBLEM] Two concrete failure cases
Modern perception pipelines typically combine a single-stage object
detector (e.g., YOLOv8~\cite{Jocher2023_yolov8}) with a multi-object
tracker (e.g., ByteTrack~\cite{Zhang2022_bytetrack}).
%
However, two systematic failure cases emerge when such pipelines operate
in motorcycle-dominant mixed traffic:

\begin{enumerate}
  \item \textbf{Track fragmentation and missed detections.}
        During dense weaving maneuvers, motorcycles occlude each other
        at irregular angles.
        The linear Kalman filter in ByteTrack predicts straight-line
        motion and loses track of vehicles executing sharp turns,
        degrading the observation quality used by downstream
        multi-agent controllers.

  \item \textbf{High false-positive accident alerts.}
        With a fixed confidence threshold, sudden braking, sharp
        lane changes, and cluster occlusions are misclassified as
        accidents, causing continuous spurious alerts to operators.
\end{enumerate}

Both failures share a common root cause: thresholds and motion models
are \emph{static}, unable to adapt to the dynamic conditions of
rush-hour density, nighttime low-contrast lighting, or wet-road glare.

% [GAP]
% TODO: reference Table~\ref{tab:related_comparison} if you add one in Sec. II
Prior systems do not bridge this gap.
Commercial platforms such as Hikvision AID~\cite{Hikvision_AID} target
highway car-only scenarios.
Academic trackers (StrongSORT, BoTrack) focus on pedestrian and vehicle
benchmarks with Kalman-based linear motion assumptions.
Context-aware detection work adjusts thresholds for weather or
illumination changes but does not consider class-conditioned motion
dynamics or dual downstream task requirements.
% TODO: expand with 2–3 more citations from Sec. II once related work is written.

% [SOLUTION — plain English]
We address these limitations with a \emph{context-aware adaptive
perception} pipeline.
A lightweight context estimator continuously reads four signals
--- vehicle density, time of day, image brightness, and recent
false-positive rate --- and uses them to adjust detection confidence
and tracking hyperparameters in real time.
A class-conditioned motion model replaces linear Kalman prediction
with the constant turn-rate and velocity (CTRV) model for motorcycles,
explicitly capturing weaving dynamics.
Finally, a dual-task track reliability score assigns each track a
quality value, enabling a single pipeline to serve both high-precision
accident alerting and high-recall signal-control observations without
task-specific retuning.

% [CONTRIBUTIONS]
\noindent\textbf{Contributions.}
The main contributions of this paper are:

\begin{itemize}
  \item \textbf{C1 — Context Estimator + Adaptive Parameter Controller:}
        a real-time module that reads four traffic context signals and
        dynamically adjusts YOLO confidence threshold, ByteTrack buffer
        length, and motion model selection, eliminating manually tuned
        fixed thresholds.

  \item \textbf{C2 — Class-Conditioned CTRV Motion Model:}
        a motorcycle-specific motion model integrated into ByteTrack
        that replaces the linear Kalman filter with CTRV prediction
        (state: $[x, y, v, \theta, \omega]^\top$), reducing track
        fragmentation during high-curvature maneuvers.

  \item \textbf{C3 — Dual-Task Track Reliability Score:}
        a lightweight scoring function that combines track age,
        detection consistency, motion smoothness, and class stability
        into a single reliability metric, with task-specific thresholds
        that balance precision for accident detection and recall for
        MAPPO-based multi-agent observations.
\end{itemize}

% [OUTLINE]
The remainder of this paper is organized as follows.
Section~\ref{sec:related} reviews related work in accident detection,
multi-object tracking, and adaptive perception.
Section~\ref{sec:method} presents the proposed pipeline in detail.
Section~\ref{sec:experiments} describes the experimental setup and
datasets.
Section~\ref{sec:results} reports ablation and comparative results.
Section~\ref{sec:conclusion} concludes with limitations and future work.
